\section{The full dyadic beta quotient as a measure algebra}
\label{sec:dyadic-beta-measure-algebra}

The integral measure describes the entire beta quotient, including its
integral lattice and arithmetic image. This also explains the successive
odd-coordinate ideals without treating their images one at a time.
Put
\[
 \mathcal U=1+4\mathbf Z_2,\qquad
 \mathcal A=\mathbf Z_2[[\mathcal U]],\qquad
 \mathfrak a=\ker(\mathcal A\longrightarrow\mathbf Z_2).
\]
We regard the completed group algebra as the integral measures on
$\mathcal U$, with convolution as multiplication and total mass as
augmentation. For $v\in\mathbf Z_2^\times$, let $\langle v\rangle$
be the unique member of $\{v,-v\}$ lying in $\mathcal U$.
The symbol $[t]$ denotes the Dirac measure at $t\in\mathcal U$.

\begin{theorem}\label{thm:full-dyadic-measure-algebra}
There is a continuous additive map
\[
 \nu:M=\ker\lambda\longrightarrow\mathcal A
\]
with the following properties:
\begin{align*}
 \nu_g(\mathcal U)&=\rho_{2,2}^I(g)/2,\\
 \int_{\mathcal U}t^{2k}\,d\nu_g(t)
   &=(1-2^{2k})\tau_{2k+1}(g)\qquad(k\ge1),\\
 \nu_{hgh^{-1}}&=\lambda_2(h)[\langle\lambda_2(h)\rangle]\nu_g.
\end{align*}
Here $\rho_{2,2}^I$ is the two-adic component of $\rho_2^I$.
Its kernel is exactly the kernel of the dyadic beta multiplier, and
\[
 \nu(M\cap H)=\nu(M)=\mathcal A,\qquad
 \nu(D\cap H)=\nu(D)=\nu(N)=\mathfrak a.
\]
Consequently the continuous additive quotients are
\[
 M/\{g\in M:B_g=1\}\simeq\mathcal A,\qquad
 D/C_\infty\simeq\mathfrak a,
\]
with the displayed conjugation action. In particular every dyadic beta
multiplier on either domain is attained on its arithmetic subgroup.
\end{theorem}

\begin{proof}
Use the integral measure $\mu_g$ with transform $q_g(X)$ from
Lemma~\ref{lem:integral-boundary-measure}, which exists on all of $M$.
Set $\nu_g=\mu_g|_{1+4\mathbf Z_2}$. The unnormalized duplication
identity, with the same conventions as
\eqref{eq:measure-duplication}, is
\[
 B_g(X,X)=X^{2\rho_{2,2}^I(g)}B_g(-1,X).
\]
Its logarithmic derivative gives
\[
 q_g(X)+q_g(-X)=2q_g(X^2)-2\rho_{2,2}^I(g).
\]
Thus, as integral measures on the additive group $\mathbf Z_2$,
\begin{equation}\label{eq:unnormalized-measure-doubling}
 \mu_g|_{2\mathbf Z_2}=[2]_*\mu_g-\rho_{2,2}^I(g)\delta_0.
\end{equation}
The total mass of $\mu_g$ is $q_g(1)=0$. Its mass on the odd
integers is therefore $\rho_{2,2}^I(g)$. Negation symmetry exchanges
the two odd residue classes modulo four, giving the claimed mass of
$\nu_g$. Integrating $t^{2k}$ in
\eqref{eq:unnormalized-measure-doubling} removes the term at zero.
The same symmetry and the definition of $\tau_{2k+1}$ give the
stated moment formula.

Additivity follows from additivity of $q_g$ on $M$, including its
measure interpretation. For $v=\lambda_2(h)$, the boundary
conjugation formula is $q_{hgh^{-1}}(X)=v q_g(X^v)$.
It sends $\mu_g$ to $v[v]_*\mu_g$. If $v\equiv1\pmod4$,
restriction to $\mathcal U$ is the pushforward by $v$; if
$v\equiv-1\pmod4$, use negation symmetry to replace this by the
pushforward by $-v$. This proves the asserted action in both cases.
All maps used are continuous in the inverse-limit topologies.

If $\nu_g=0$, its total mass gives $\rho_{2,2}^I(g)=0$.
Symmetry makes the whole odd restriction of $\mu_g$ zero, and
\eqref{eq:unnormalized-measure-doubling} gives
$\mu_g=[2]_*\mu_g$. Iteration shows that $\mu_g$ is supported on
$\bigcap_n2^n\mathbf Z_2=\{0\}$. Its total mass is zero, so it is
the zero measure. Hence $q_g=0$. The beta differential identity
$D_Y\log B_g(X,Y)=q_g(XY)-q_g(Y)$ and $B_g(X,1)=1$ give $B_g=1$.
The converse is immediate. This proves the kernel assertion; it
concerns the dyadic component and does not determine the other
components of $\rho_2^I$.

Choose $\sigma\in H\cap M$ with $\rho_2^I(\sigma)=2$ and all
odd-prime rational Kummer coordinates zero, and $h\in H$ with
$t=\lambda(h)$ and $t_2=3$, as in the earlier arithmetic witnesses.
Write $\gamma=[-3]\in\mathcal A$. The element $-3$ topologically
generates $\mathcal U$: one has
$v_2((-3)^{2^n}-1)=n+2$, so its image generates every quotient
$(1+4\mathbf Z_2)/(1+2^{n+2}\mathbf Z_2)$. Therefore
\[
 \mathcal A\simeq\mathbf Z_2[[Z]],\qquad \gamma=1+Z,
 \qquad \mathfrak a=Z\mathcal A.
\]
The measure $\nu_\sigma$ has augmentation one, so it is a unit of
this power-series ring. The compact additive subgroup
$\nu(H\cap M)$ is a closed $\mathbf Z_2$-submodule, and
conjugation by $h$ acts as $3\gamma$. Multiplication by three is
invertible on a closed $\mathbf Z_2$-submodule. The image therefore
contains every polynomial in $\gamma$ times $\nu_\sigma$, and by
closure every element of $\mathcal A\nu_\sigma=\mathcal A$.
This proves the full arithmetic image.

Finally let $z=h\sigma h^{-1}\sigma^{-t}$. The full profinite
Kummer laws put $z$ in $D\cap H$, and
\[
 \nu_z=3(\gamma-1)\nu_\sigma.
\]
This generates the ideal $\mathfrak a$ as a closed
$\mathcal A$-module. Normality of $D$ and the same conjugation
argument show that $\nu(D\cap H)$ contains this ideal.
Conversely $D\subseteq N$, and the mass formula puts $\nu(N)$
in $\mathfrak a$. All the claimed image equalities follow.
The kernel on $D$ is $C_\infty$ by its beta characterization.
A continuous surjection from a compact group to a Hausdorff group
induces the claimed topological quotient isomorphisms.
\end{proof}

\begin{corollary}\label{cor:unnormalized-odd-images}
Put $T_0=M$ and
\[
 T_r=M\cap\bigcap_{j=1}^r\ker\tau_{2j+1}\qquad(r\ge1).
\]
For every $r\ge1$, the raw coefficient
$m_{2r+1}=[S^{2r}T]f$ is additive on $T_{r-1}$, has conjugation
weight $\lambda_2^{2r+1}$, and has exact full and arithmetic image
\[
 m_{2r+1}(T_{r-1})=m_{2r+1}(T_{r-1}\cap H)
                  =2^{2r-2-v_2(r)}\mathbf Z_2.
\]
Its kernel is $T_r$. More precisely, in the preceding identification
of $\mathcal A$ one has
\[
 \nu(T_r)=\nu(T_r\cap H)
       =\prod_{j=1}^r(\gamma-9^j)\mathcal A,
 \qquad
 \nu(C_r)=\nu(C_r\cap H)
       =\prod_{j=0}^r(\gamma-9^j)\mathcal A.
\]
\end{corollary}

\begin{proof}
The moment $\int t^{2j}\,d\nu(t)$ is the evaluation homomorphism
$\operatorname{ev}_{9^j}:\mathcal A\to\mathbf Z_2$ that sends
$\gamma$ to $9^j$. Its scalar factor $1-2^{2j}$ in the theorem
is a unit. For $\alpha\in2\mathbf Z_2$, ordinary power-series
division gives
\[
 \ker\bigl(\mathbf Z_2[[Z]]\xrightarrow{Z\mapsto\alpha}
                  \mathbf Z_2\bigr)=(Z-\alpha).
\]
Indeed the coefficients of $(F(Z)-F(\alpha))/(Z-\alpha)$ are
convergent sums of integral coefficients times powers of $\alpha$.
Applying this division successively proves that the intersection
of the evaluation kernels at the distinct points $9^j$ is the
product of the displayed linear factors. At each step the
previous factors have nonzero value in $\mathbf Z_2$ at the new
point, so cancellation is valid even though that value need not
be a unit. Theorem~\ref{thm:full-dyadic-measure-algebra} lifts
every measure in either resulting ideal to the corresponding
arithmetic subgroup, proving the two exact ideal formulas.

On $T_{r-1}$, reflection and vanishing of the lower moments make
$q_g(1+S)$ start in degree $2r$, so
\[
 m_{2r+1}=\frac{2}{(2r)!}\tau_{2r+1}.
\]
Evaluation of the ideal for $T_{r-1}$ at $\gamma=9^r$ has image
\[
 \prod_{j=1}^{r-1}(9^r-9^j)\mathbf Z_2.
\]
Its valuation is $3(r-1)+v_2((r-1)!)$, by
$v_2(9^k-1)=3+v_2(k)$. After the factor $2/(2r)!$ the valuation is
\[
 1+3(r-1)+v_2((r-1)!)-r-v_2(r!)=2r-2-v_2(r).
\]
The empty product for $r=1$ gives the same formula. The exact
images and kernel follow. Additivity and the weight law follow
from the corresponding properties of the moments.
\end{proof}

\begin{corollary}\label{cor:beta-quadratic-invariants}
Among continuous invariant quadratic characters which factor through
the dyadic beta multiplier, the space on $M$ is one-dimensional,
generated by $\rho_{2,2}^I/2\pmod2$. On each of $N$ and $D$ it
is one-dimensional, generated by $\eta_3=m_3/8\pmod2$.
\end{corollary}
\begin{proof}
Invariance under the chosen $h$ forces a quadratic character of
$\mathcal A$ to kill $2\mathcal A$ and $(3\gamma-1)\mathcal A$.
Modulo two the latter is $(\gamma-1)\mathcal A$, so the invariant
quotient is $\mathcal A/(2,\gamma-1)\simeq\mathbf F_2$.
Augmentation modulo two is invariant under every multiplier by
the action in the theorem and is nonzero.
For $\mathfrak a=(\gamma-1)\mathcal A$, the analogous quotient
is $\mathfrak a/(2\mathfrak a+\mathfrak a^2)\simeq\mathbf F_2$.
Write $\nu=(\gamma-1)A$. Its second moment is $8A(9)$, whence
\[
 m_3=\tau_3=-\frac{8}{3}A(9),\qquad
 \eta_3=A(1)\pmod2.
\]
This character is invariant and nonzero by the arithmetic image.
The claimed descriptions follow from the exact beta quotients.
\end{proof}
